% Options for packages loaded elsewhere
\PassOptionsToPackage{unicode}{hyperref}
\PassOptionsToPackage{hyphens}{url}
\PassOptionsToPackage{dvipsnames,svgnames*,x11names*}{xcolor}
%
\documentclass[
  11pt,
]{article}
\usepackage{amsmath,amssymb}
\usepackage{lmodern}
\usepackage{iftex}
\ifPDFTeX
  \usepackage[T1]{fontenc}
  \usepackage[utf8]{inputenc}
  \usepackage{textcomp} % provide euro and other symbols
\else % if luatex or xetex
  \usepackage{unicode-math}
  \defaultfontfeatures{Scale=MatchLowercase}
  \defaultfontfeatures[\rmfamily]{Ligatures=TeX,Scale=1}
  \setmainfont[]{Helvetica Neue}
  \setsansfont[]{Helvetica Neue}
  \setmathfont[]{STIX Two Math}
\fi
% Use upquote if available, for straight quotes in verbatim environments
\IfFileExists{upquote.sty}{\usepackage{upquote}}{}
\IfFileExists{microtype.sty}{% use microtype if available
  \usepackage[]{microtype}
  \UseMicrotypeSet[protrusion]{basicmath} % disable protrusion for tt fonts
}{}
\makeatletter
\@ifundefined{KOMAClassName}{% if non-KOMA class
  \IfFileExists{parskip.sty}{%
    \usepackage{parskip}
  }{% else
    \setlength{\parindent}{0pt}
    \setlength{\parskip}{6pt plus 2pt minus 1pt}}
}{% if KOMA class
  \KOMAoptions{parskip=half}}
\makeatother
\usepackage{xcolor}
\IfFileExists{xurl.sty}{\usepackage{xurl}}{} % add URL line breaks if available
\IfFileExists{bookmark.sty}{\usepackage{bookmark}}{\usepackage{hyperref}}
\hypersetup{
  colorlinks=true,
  linkcolor={accent},
  filecolor={Maroon},
  citecolor={Blue},
  urlcolor={Blue},
  pdfcreator={LaTeX via pandoc}}
\urlstyle{same} % disable monospaced font for URLs
\usepackage[margin=18mm]{geometry}
\setlength{\emergencystretch}{3em} % prevent overfull lines
\providecommand{\tightlist}{%
  \setlength{\itemsep}{0pt}\setlength{\parskip}{0pt}}
\setcounter{secnumdepth}{-\maxdimen} % remove section numbering
\usepackage{xcolor}
\usepackage{titlesec}
\definecolor{textmain}{HTML}{1F2328}
\definecolor{muted}{HTML}{6B7280}
\definecolor{accent}{HTML}{2563EB}
\definecolor{linegray}{HTML}{D9DEE7}
\color{textmain}
\renewcommand{\familydefault}{\sfdefault}
\setlength{\parindent}{0pt}
\setlength{\parskip}{5pt}
\setlength{\emergencystretch}{3em}
\titleformat{\section}{\fontsize{16}{20}\selectfont\bfseries\color{textmain}}{\thesection}{8pt}{}
\titleformat{\subsection}{\fontsize{13}{16}\selectfont\bfseries\color{accent}}{\thesubsection}{6pt}{}
\titlespacing*{\section}{0pt}{18pt}{6pt}
\titlespacing*{\subsection}{0pt}{12pt}{4pt}
\ifLuaTeX
  \usepackage{selnolig}  % disable illegal ligatures
\fi

\title{Lecture 6 Slides ↔ Matousek Paper Correspondence}
\usepackage{etoolbox}
\makeatletter
\providecommand{\subtitle}[1]{% add subtitle to \maketitle
  \apptocmd{\@title}{\par {\large #1 \par}}{}{}
}
\makeatother
\subtitle{matousek-partition-tree / docs}
\author{}
\date{}

\begin{document}
\maketitle

{
\hypersetup{linkcolor=}
\setcounter{tocdepth}{3}
\tableofcontents
}
\hypertarget{lecture-06-slides-85-95-matousek-paper-calculation-extraction}{%
\section{Lecture 06 Slides 85-95: Matousek Paper Calculation
Extraction}\label{lecture-06-slides-85-95-matousek-paper-calculation-extraction}}

Source paper: Jiri Matousek, ``Efficient Partition Trees'', Discrete \&
Computational Geometry 8:315-334, 1992.

Scope: only the paper content corresponding to the lecture slides on
partition trees and simplicial partitions.

\hypertarget{background-warmup-four-way-partition}{%
\subsection{1. Background Warmup: Four-Way
Partition}\label{background-warmup-four-way-partition}}

This part is not the main Matousek theorem. It is the older/simple
partition-tree idea used in the slides before introducing simplicial
partitions.

The plane is recursively divided into four regions, each containing:

\[
\frac n4
\]

points.

For a halfplane query, the query boundary is a line. In the worst case,
this line intersects three of the four regions, so the query recurses
into three children.

Thus the query recurrence is:

\[
T(n)=3T(n/4)+O(1)
\]

Solving:

\[
T(n)=O(n^{\log_4 3})
\]

Since:

\[
\log_4 3\approx 0.792
\]

we get:

\[
T(n)=O(n^{0.792})
\]

This is the bound shown on slide 89.

\hypertarget{paper-definition-simplicial-partition}{%
\subsection{2. Paper Definition: Simplicial
Partition}\label{paper-definition-simplicial-partition}}

The paper defines a simplicial partition of P as:

\[
\Pi=\lbrace (P_1,\Delta_1),\ldots,(P_m,\Delta_m) \rbrace
\]

where:

\begin{itemize}
\tightlist
\item
  the sets \(P_i\) are pairwise disjoint,
\item
  together they form a partition of P,
\item
  each \(\Delta_i\) is a relatively open simplex containing \(P_i\).
\end{itemize}

In dimension 2, a simplex is a triangle, so the slide writes:

\[
F(P)=\lbrace (P_1,t_1),\ldots,(P_r,t_r) \rbrace
\]

where each \(t_i\) is a triangle containing \(P_i\).

The paper allows lower-dimensional simplices for degenerate point sets,
but the slides only use the planar triangle picture.

\hypertarget{paper-definition-crossing-number}{%
\subsection{3. Paper Definition: Crossing
Number}\label{paper-definition-crossing-number}}

The paper says a hyperplane h crosses a simplex \(\Delta\) if:

\[
h\cap\Delta\ne\emptyset
\]

and:

\[
\Delta\not\subset h
\]

In the plane, hyperplanes are lines.

So the crossing number of a line h relative to \(\Pi\) is:

\[
\#\lbrace \Delta_i:h\text{ crosses }\Delta_i \rbrace
\]

The crossing number of the whole partition is:

\[
\max_h \#\lbrace \Delta_i:h\text{ crosses }\Delta_i \rbrace
\]

This is exactly slide 91:

\[
\text{crossing number}=\text{maximal number of triangles intersected by a line}
\]

\hypertarget{paper-class-size-condition-and-slide-fine-condition}{%
\subsection{4. Paper Class-Size Condition and Slide ``Fine''
Condition}\label{paper-class-size-condition-and-slide-fine-condition}}

The paper uses parameter s and requires:

\[
s\le |P_i|<2s
\]

It also sets:

\[
r=\frac ns
\]

Therefore:

\[
2s=\frac{2n}{r}
\]

So the paper condition implies:

\[
|P_i|< \frac{2n}{r}
\]

The slide calls a simplicial r-partition fine if:

\[
|P_i|\le \frac{2n}{r}
\]

So the slide's fine condition is the upper-bound part of the paper's
stronger condition.

\hypertarget{paper-theorem-3.1-partition-theorem}{%
\subsection{5. Paper Theorem 3.1: Partition
Theorem}\label{paper-theorem-3.1-partition-theorem}}

The paper states in fixed dimension d:

Given P, n, and s with:

\[
2\le s< n
\]

set:

\[
r=\frac ns
\]

Then there exists a simplicial partition \(\Pi\) such that:

\[
s\le |P_i|<2s
\]

and the crossing number is:

\[
O(r^{1-1/d})
\]

For the lecture slides, d equals 2. Therefore:

\[
r^{1-1/d}=r^{1-1/2}=r^{1/2}=\sqrt r
\]

So in the plane:

\[
\text{crossing number}=O(\sqrt r)
\]

This is the theorem stated on slide 92:

\[
\text{fine simplicial partition of size }r \text{ and crossing number }O(\sqrt r)
\]

\hypertarget{paper-construction-time-used-by-slides}{%
\subsection{6. Paper Construction Time Used by
Slides}\label{paper-construction-time-used-by-slides}}

The paper's exact algorithmic statement is in Theorem 4.7.

The slide uses the clean version:

\[
O(n^{1+\varepsilon})
\]

construction time for a planar fine simplicial partition with optimal
crossing number:

\[
O(\sqrt r)
\]

This corresponds to Theorem 4.7(iii), which says for any s, a simplicial
partition with crossing number:

\[
O(r^{1-1/d})
\]

can be constructed in:

\[
O(n^{1+\delta})
\]

Renaming the paper's fixed positive constant \(\delta\) as the lecture's
\(\varepsilon\), in 2D this becomes:

\[
O(n^{1+\varepsilon})
\]

preprocessing/construction time.

\hypertarget{building-the-partition-tree-from-the-theorem}{%
\subsection{7. Building the Partition Tree from the
Theorem}\label{building-the-partition-tree-from-the-theorem}}

At a node containing n points:

\begin{itemize}
\tightlist
\item
  build a fine simplicial r-partition,
\item
  create r children,
\item
  each child contains at most:
\end{itemize}

\[
\frac{2n}{r}
\]

points, - store the triangle/simplex for each child.

For a halfplane query, the boundary is a line.

At the current node:

\begin{enumerate}
\def\labelenumi{\arabic{enumi}.}
\tightlist
\item
  Check all r child triangles.
\item
  If a child triangle is fully inside the query halfplane, add its
  stored aggregate answer.
\item
  If it is disjoint, ignore it.
\item
  If it is crossed by the query boundary line, recurse into that child.
\end{enumerate}

Because the crossing number is:

\[
O(\sqrt r)
\]

the query recurses into only:

\[
O(\sqrt r)
\]

children.

This gives the slide recurrence:

\[
T(n)=r+\sqrt r\cdot T(2n/r)
\]

The first term r is the cost of scanning all r children at the current
node.

The recursive term appears because:

\begin{itemize}
\tightlist
\item
  at most \(\sqrt r\) children are crossed,
\item
  each crossed child has at most \(2n/r\) points.
\end{itemize}

\hypertarget{solving-the-slide-query-recurrence}{%
\subsection{8. Solving the Slide Query
Recurrence}\label{solving-the-slide-query-recurrence}}

The recurrence is:

\[
T(n)=r+\sqrt r\cdot T(2n/r)
\]

The slides choose r as a sufficiently large constant depending only on
\(\varepsilon\). Therefore r does not grow with n.

We prove:

\[
T(n)=O(n^{1/2+\varepsilon})
\]

Assume inductively:

\[
T(m)\le C m^{1/2+\varepsilon}
\]

Then:

\[
T(n)\le r+\sqrt r\cdot C(2n/r)^{1/2+\varepsilon}
\]

Simplify the recursive term:

\[
\sqrt r\cdot (2n/r)^{1/2+\varepsilon} = r^{1/2}\cdot 2^{1/2+\varepsilon}n^{1/2+\varepsilon}r^{-1/2-\varepsilon}
\]

Cancel the powers of r:

\[
r^{1/2}r^{-1/2-\varepsilon}=r^{-\varepsilon}
\]

So:

\[
\sqrt r\cdot (2n/r)^{1/2+\varepsilon} = 2^{1/2+\varepsilon}r^{-\varepsilon}n^{1/2+\varepsilon}
\]

Thus:

\[
T(n)\le r+C2^{1/2+\varepsilon}r^{-\varepsilon}n^{1/2+\varepsilon}
\]

To make the induction close, choose r so large that:

\[
2^{1/2+\varepsilon}r^{-\varepsilon}<1
\]

Equivalently:

\[
r>2^{(1/2+\varepsilon)/\varepsilon}
\]

The slide gives one convenient constant choice of this type:

\[
r=2\cdot 2^{1/(2\varepsilon)}
\]

Up to constant-factor slack, this is just saying:

\[
r=\Theta(2^{1/(2\varepsilon)})
\]

With such a constant r:

\[
T(n)=O(n^{1/2+\varepsilon})
\]

This is slide 95.

\hypertarget{space-bound-in-the-slides}{%
\subsection{9. Space Bound in the
Slides}\label{space-bound-in-the-slides}}

At a node with m points, the partition creates r children.

Because r is chosen as a constant depending only on \(\varepsilon\),
each node has constant branching factor.

Each point belongs to exactly one child at each level.

The partition tree stores:

\begin{itemize}
\tightlist
\item
  one aggregate value per child,
\item
  one simplex/triangle per child,
\item
  leaves for constant-size point sets.
\end{itemize}

The total number of tree nodes is linear in the number of leaves, and
the leaves partition the original point set.

Therefore total space is:

\[
O(n)
\]

This is slide 93.

\hypertarget{preprocessing-bound-in-the-slides}{%
\subsection{10. Preprocessing Bound in the
Slides}\label{preprocessing-bound-in-the-slides}}

At each node, the expensive step is constructing the simplicial
partition.

The lecture uses Matousek's construction-time guarantee:

\[
O(m^{1+\varepsilon})
\]

for a node with m points, with r treated as a constant depending on
\(\varepsilon\).

Since the children form a partition of the node's points and each child
has at most a constant fraction of the parent points, the total over the
recursive tree remains:

\[
O(n^{1+\varepsilon})
\]

Intuition:

At one level with subproblem sizes \(m_1,\ldots,m_k\), where:

\[
\sum_j m_j=n
\]

and each \(m_j\) is a constant fraction of its parent, the superlinear
exponent makes the geometric decay dominate across levels.

Hence total preprocessing:

\[
O(n^{1+\varepsilon})
\]

This is slide 94.

\hypertarget{final-slide-theorem}{%
\subsection{11. Final Slide Theorem}\label{final-slide-theorem}}

Combining:

\begin{itemize}
\tightlist
\item
  Matousek 2D simplicial partition theorem:
\end{itemize}

\[
O(\sqrt r)
\]

crossing number,

\begin{itemize}
\tightlist
\item
  constant-r recursive partition tree,
\item
  recurrence:
\end{itemize}

\[
T(n)=r+\sqrt r\cdot T(2n/r)
\]

\begin{itemize}
\tightlist
\item
  sufficiently large constant r depending on \(\varepsilon\),
\end{itemize}

gives:

\[
T(n)=O(n^{1/2+\varepsilon})
\]

The resulting planar halfplane range searching structure has:

\[
O(n)
\]

space,

\[
O(n^{1+\varepsilon})
\]

construction time,

and:

\[
O(n^{1/2+\varepsilon})
\]

query time.

This is exactly slide 96.

\end{document}
